% tps/tp7.tex : Déploiement et benchmark sur l'ESP32-S3
% Squelette : questions écrites, prose et code à faire.
% Décision ouverte : framework embarqué (ESP-DL ou TFLite Micro).

\chapter{Déploiement et benchmark sur l'ESP32-S3}

\section*{Motivation}

Le modèle compressé des TP 5 et 6 doit maintenant tourner sur la carte, sur des images de la caméra, dans les limites de mémoire et de temps de l'ESP32-S3. Ce TP convertit le modèle pour le moteur d'inférence embarqué, l'intègre dans un sketch Arduino-ESP32, et mesure ce qui compte pour l'application : latence, mémoire occupée, précision sur la carte. \\

% TODO prose : moteur d'inférence retenu, utilisé en boîte noire ; chaîne image caméra -> prétraitement -> inférence -> classe.

\begin{figure}[H]
\centering
\resizebox{0.8\linewidth}{!}{% diagrams/tp7/deployment.tex : du modèle PyTorch à la classe affichée par la carte
\begin{tikzpicture}[
    font=\sffamily\small, >=Stealth,
    pc/.style={draw=ECE, thick, fill=ECE!10, rounded corners=2pt, minimum height=1cm, minimum width=2.2cm, align=center, inner sep=4pt},
    mcu/.style={pc, draw=gamboge, fill=gamboge!15},
    lbl/.style={font=\sffamily\scriptsize, text=black!70, align=center}
]
\node[pc] (pt) {modèle PyTorch\\TP 5 et 6};
\node[pc, right=0.9cm of pt] (onnx) {fichier ONNX\\vérifié (TP 6)};
\node[pc, right=0.9cm of onnx] (conv) {outil de conversion\\du framework embarqué};
\node[mcu, right=0.9cm of conv] (bin) {modèle embarqué\\en flash};
\draw[->, thick] (pt) -- (onnx);
\draw[->, thick] (onnx) -- (conv);
\draw[->, thick] (conv) -- (bin);
\node[mcu, below=1.1cm of pt] (cam) {caméra\\OV3660};
\node[mcu, right=0.9cm of cam] (pre) {prétraitement\\redimension, normalisation};
\node[mcu, right=0.9cm of pre] (inf) {inférence};
\node[mcu, right=0.9cm of inf] (cls) {classe\\série ou LED RGB};
\draw[->, thick] (cam) -- (pre);
\draw[->, thick] (pre) -- (inf);
\draw[->, thick] (inf) -- (cls);
\draw[->, thick, dashed] (bin.south) -- ++(0,-0.4) -| (inf.north);
\node[lbl, anchor=east] at ([xshift=-0.2cm]pt.west) {PC};
\node[lbl, anchor=east] at ([xshift=-0.2cm]cam.west) {carte};
\end{tikzpicture}
}
\caption{Chaîne de déploiement : conversion sur PC, puis capture, prétraitement et inférence sur la carte.}
\label{fig:tp7-deploy}
\end{figure}

\section{Conversion et intégration}

\EQ{Convertir le modèle ONNX du TP 6 avec l'outil du framework embarqué. Relever la taille du modèle converti et la comparer à celle du modèle int8 du TP 5.}

\EQ{Intégrer le modèle dans le sketch fourni. Relever la mémoire occupée en flash, en SRAM et en PSRAM après compilation et au démarrage.}

\TQ{Où placer les poids du modèle (flash, SRAM, PSRAM) et la mémoire de travail des activations ? Quel compromis entre capacité et vitesse d'accès ?}

\begin{imp}
\textbf{À retenir :} vérifier la taille du modèle converti et la mémoire libre avant d'écrire le reste du sketch. Un modèle qui ne tient pas en SRAM passe en PSRAM, plus lente : l'emplacement des poids et des activations fixe une partie de la latence.
\end{imp}

\section{Validation sur la carte}

\EQ{Envoyer à la carte, par le port série, 50 images du jeu de test. Comparer les classes prédites par la carte à celles du modèle sur PC. Combien de prédictions diffèrent ?}

\EQ{Mesurer avec \texttt{esp\_timer\_get\_time()} la durée du prétraitement de l'image et la durée de l'inférence, moyennées sur 100 images. Combien d'images par seconde la carte peut-elle classer ?}

\EQ{Comparer latence et précision sur la carte pour le modèle \texttt{float32}, le modèle int8 et le modèle élagué et quantifié.}

\EQ{Boucler capture, prétraitement, inférence et affichage de la classe sur le port série (ou par la LED RGB). Présenter devant la caméra des images imprimées des 4 classes et relever le taux de bonnes réponses sur 20 essais par classe. Comparer à la précision de test mesurée sur PC.}

\begin{imp}
\textbf{À retenir :} un modèle embarqué se valide sur la carte, avec des images de la caméra. La précision sur PC est une borne supérieure : prétraitement différent, éclairage et quantification la dégradent. Mesurer la latence de toute la chaîne, pas seulement celle de l'inférence.
\end{imp}
